\documentclass[journal]{IEEEtran}

\usepackage{cite}
\usepackage{amsmath,amssymb,amsfonts}
\usepackage{graphicx}
\usepackage{textcomp}
\usepackage{xcolor}
\usepackage{listings}
\usepackage[inline, shortlabels]{enumitem}

\lstset{
	basicstyle=\footnotesize\ttfamily,
	breaklines=true,
	frame=single,
	captionpos=b,
	tabsize=2,
	showstringspaces=false,
	commentstyle=\color{gray},
	keywordstyle=\color{blue},
	stringstyle=\color{red},
	frame=none
}

\begin{document}
	
	\title{Debugging Issues in the Removals System}
	
	\author{Mindaugas Taujanskas}
	
	\maketitle
	
	\section{Creating a Business as a Customer}
	
	During the creation of tests for the business creation logic, a bug was discovered regarding the permissibility of customers creating business resources, contrary to the business logic that indicates only service providers may register businesses. The test case is as follows:
	
	\begin{lstlisting}[language=Python]
def test_appropriate_role_permission(
	db_guest_cursor,
	with_valid_user
) -> None:
	*_, session = with_valid_user("customer")
	fetch_result = db.proc_create_business(
		session.data.token,
		business_name="Sample business name",
		vat_no=gen_rand_alnum(length=11),
		crn_no=gen_rand_alnum(length=8),
		utr_no=gen_rand_alnum(length=10),
		num_employees=2
	)
	assert not fetch_result.success
	\end{lstlisting}
	
	The case fails on \lstinline[language=Python]|assert not fetch_result.success|, indicating that the database registered the business resource despite the registering user being a customer. To ensure the database handles this validation, check the procedural definition for the corresponding function:
	
	\begin{lstlisting}[language=SQL]
CREATE FUNCTION create_business(
	p_token TEXT,
	...
)
RETURNS result_t AS $$
DECLARE
	user_session JSON;
	stored_business_id INTEGER;
BEGIN
	user_session := decode_token(p_token);
	
	IF user_session IS NULL THEN
		RETURN make_error_result(
			'INVALID_SESSION',
			'Invalid session token'
		);
	ELSIF user_session->>'user_role' <> 'service-provider' THEN
		RETURN make_error_result(
			'INSUFFICIENT_PERMISSIONS',
			'Invalid role for creating business'
		);
	END IF;
	
	...
	
	RETURN make_success_result();
END;
$$ LANGUAGE plpgsql SECURITY DEFINER;
	\end{lstlisting}
	
	As read, this appears to appropriately check the session data for the user role and is thus non-problematic, but upon further evaluation \lstinline[language=SQL]|user_session->>'user_role'| contains a subtle typo, and must be \lstinline[language=SQL]|user_session->>'role'|, as defined by the registration procedures:
	
	\begin{lstlisting}[language=SQL]
token := sign(
	json_build_object(
		'user_id', new_user_id,
		'email', p_email,
		'role', p_user_role  -- <<< HERE
	),
	utils.get_app_config_value('jwt_secret'),
'HS256'
);
	\end{lstlisting}
	
	After correcting the typo, the test passes, meaning our business logic is correctly implemented at the database-level.
\end{document}